\scanpage{518}
\mspage{429}
% B01
Pour prouver que $x=\underset I{\unsure{\underline{\lim}|\varinjlim}}\,x_i$ existe dans $E$, et commute à $f$, on peut remplacer $I$ par une sous-catégorie cofinale. Cela nous ramène, grâce au lemme, au cas où $(s_i)_{i\in I}$ est constant, donc au cas où le système inductif provient d'un système inductif dans une catégorie fibre \inserted{$E_s$}. Mais par a), on sait que sa limite $x$ dans $E_s$ est une limite dans $E$, \unsure{d'où|\deleted{d'où}} l'existence de cette dernière, et la commutation à $f$, puisque $\underset i{\underline{\lim}}\,\deleted{\unsure{f}}(s_i)=s$ ($\underline{\lim}$ d'un syst. inductif filtrant constant).

% B02
c) Il faut prouver que si $x\in\mathrm{Ob}\,E_s$ est $\pi$-accessible \deleted{\unsure{il est}} dans $E_s$, il l'est dans $E$, i.e. que pour tt syst. inductif $(y_i)_{i\in I}$ grd devant $\pi$, la flèche canonique
% B03
{\csname tagsleft@true\endcsname
\[\varinjlim_i\mathrm{Hom}(x,y_i)\longrightarrow\mathrm{Hom}(x,y)\tag*{(**)}\]}
% B04
est bijective. Or \unsure{comme|\deleted{le} comme} \inserted{par $2^\circ$)} $\pi\geq\pi_1\geq\mathrm{Sup}(\mathrm{card}\,\unsure{E|C},\underset{\lambda}{\mathrm{Sup}}\pi_\lambda)$, il résulte encore du lemme que, quitte à remplacer $I$ par une ss-cat. cofinale, \unsure{OPS|OP S} que le système inductif $(y_i)_{i\in I}$ provient d'un système dans \deleted{\unsure{E}} une catégorie fibre $E_t$, et que $y$ est sa $\unsure{\underset{\longrightarrow}{\mathrm{limite}}|\varinjlim^{i}|\varinjlim}$ dans $E_t$.

% B05
\deleted{\unsure{On a vu de}} \inserted{Mais alors} les ens $\mathrm{Hom}(x,y)$ et $\mathrm{Hom}(x,y_i)$, et les \deleted{\unsure{applications}} applications de transition entre celles-ci, sont encore dans $\unsure{\overline{\mathrm E}\mathrm{ns}|\mathrm{Ens}}_{/\mathrm{Hom}(s,t)}$, et on est ramené

% archive sheet number: 517
